Generative AI is both the subject of this project and a tool used during it, so
the distinction between the two needs stating plainly.

\textbf{As the object of study.} The artefact evaluated in this report is itself a
multi-agent large language model system. All agents run on Llama 3.1 (8B) served
locally through Ollama, at the temperatures recorded in
Table~\ref{tab:agent-config}. Every prompt is reproduced or summarised in
Appendix~\ref{app:prompts}, and every negotiation transcript cited in
Chapter~\ref{ch:evaluation} is committed to the project repository. The details
needed to reproduce any figure in this report are given in
Appendix~\ref{app:repro}.

\textbf{As a tool during the work.} I used Claude (Anthropic) as a coding and
research assistant throughout development, in the same way a developer uses an IDE
assistant: drafting and refactoring Python, writing test cases, searching case law
databases for candidate judgments, and reviewing my own prose for clarity. Three
constraints governed that use, and they are worth recording because they shaped
the results. First, no case was written into the evaluation corpus on the
assistant's say-so; every citation, contract value, and disposition in
Table~\ref{tab:corpus} was checked against BAILII, the National Archives case law
service, or a published law-firm case note before being committed, and seven
plausible-looking candidates were rejected during that check. Second, no
experimental number in this report was produced by an assistant; all figures come
from analysis scripts in the repository, run against committed logs, and the
headline figures were re-run immediately before submission. Third, the
interpretation, the argument, and the decisions about what counts as an honest
claim are mine.

The irony is not lost on me that a project whose central finding concerns model
fabrication was built with substantial model assistance. That is partly why the
verification discipline described in Section~\ref{sec:verification-discipline} was
applied to the write-up as well as to the artefact: at least two errors in earlier
drafts of the evaluation documents were caught by checking a claim against the data
it purported to describe, rather than against the plausibility of the sentence.
